\subsection{中值定理}

2.2.1. 设 \(x,y\) 属于 \(E\)，\([x,y]\) 是连接这两点的闭线段。此外，设 \(f\) 是从 \([x,y]\) 的某邻域到空间 \(F\) 的映射，并在 \([x,y]\) 的每一点处可微。则 \(f(x)-f(y)\) 属于当 \(Df(z).(x-y)\) 遍历 \(z\) 时的点集 \([x,y]\) 的闭凸包。

2.2.2. 设 U 是 E 的连通开子集，f 是从 U 到
F 的映射，并在 U 的每一点处导数为零；则 f 在 U 上为常值。

2.2.3. 设 U 是 E 的凸开子集，f 是从 U 到 F 的映射，并在 U 的每一点处可微。给定 F 上的连续半范数 γ 和实数 M ≥ 0，下列条件等价：

\begin{enumerate}
\def\labelenumi{(\roman{enumi})}
\item
  对每个 \(x\)  \(\mathbf{U}\)，都有 \(\| Df(x)\|_{\gamma}\leqslant M\)。
\item
  对 U 中任意 \(x\) 和任意 \(y\)\(U\)，都有 \(\| f(x) - f(y)\|_{\gamma} \leqslant M . \| x - y\| \)。
\end{enumerate}

2.2.4. 设 U 是 E 中一点 \(x_0\) 的一个邻域，\(f\) 是在 U 去掉 \(x_0\) 的补集上定义且取值于 F 的函数。假设 \(f\) 在 U 中每个满足 \(Df(x)\) 的点 \(x\) 处都有导数 \(x \neq x_0\)，且函数 \(x \mapsto Df(x)\) 在 \(D_0\) 趋于 \(x\) 时有极限 \(x_0\)。那么，当  (E) \(\geqslant 2\) 时，\(f\) 在 \(x_0\) 处有极限，并且连续延拓到整个 U 的函数 \(f\) 在 x＿0 处可微、导数为 \(D_0\)；若 \(x_0\) 且假设 \(f\) 在 \(x_0\) 处有极限，同样结论也成立。

